\documentclass[11pt]{article}
\usepackage[margin=1in]{geometry}
\usepackage{parskip}
\usepackage{hyperref}

\title{Lecture Notes Script (Slide-by-Slide): Lecture 3}
\author{Hui-Jun Chen}
\date{\today}

\begin{document}
\maketitle

\section*{Lecture 3 (AD--AS: add inflation + short-run vs medium-run)}

\begin{enumerate}
\item \textbf{Why we need AD--AS after IS--LM}
\begin{itemize}
\item Lecture 2 gave AD: mapping $P\mapsto Y$.
\item Inflation is about how $P$ moves, so we add AS / price adjustment.
\end{itemize}

\item \textbf{Learning objectives}
\begin{itemize}
\item Derive AD and shifts; define FE/CAS; combine to get AD--AS equilibrium.
\item Predict $(Y,P)$ for demand vs supply shocks, short-run vs medium-run.
\end{itemize}

\item \textbf{Roadmap: what each set of figures is doing}
\begin{itemize}
\item AD from IS + policy.
\item AS from RBC supply benchmark (FE).
\item AD--AS comparative statics: demand vs supply.
\end{itemize}

\item \textbf{From IS--LM to IS--MP (modern central banks)}
\begin{itemize}
\item Modern CB targets interest rate; represent policy with an MP rule.
\item Keep demand logic but match institutions.
\end{itemize}

\item \textbf{Step 1: IS as goods-demand relation}
\begin{itemize}
\item IS: lower real rate raises demand $\Rightarrow$ higher $Y$ in equilibrium.
\item Later: this becomes Euler/NK IS.
\end{itemize}

\item \textbf{Step 2: IS + MP implies AD}
\begin{itemize}
\item MP pins down rate; IS pins down $Y$ given rate.
\item For each $P$, you get an implied $Y$.
\end{itemize}

\item \textbf{AD curve (IS--MP--AD)}
\begin{itemize}
\item AD is $(P,Y)$ pairs consistent with goods-market equilibrium + policy.
\item Lower $P$ typically loosens conditions $\Rightarrow$ higher $Y$.
\end{itemize}

\item \textbf{AD shifts: fiscal expansion ($G\uparrow$)}
\begin{itemize}
\item $G\uparrow$ raises demand directly $\Rightarrow$ AD shifts right.
\end{itemize}

\item \textbf{AD shifts: easier monetary policy}
\begin{itemize}
\item Easier policy lowers rates $\Rightarrow$ higher demand $\Rightarrow$ AD right.
\end{itemize}

\item \textbf{Aggregate Supply: FE benchmark}
\begin{itemize}
\item FE/CAS from RBC supply block: long-run anchor for output.
\item Short run can deviate; medium run returns toward FE as prices adjust.
\end{itemize}

\item \textbf{CAS shifts: productivity falls ($A\downarrow$)}
\begin{itemize}
\item Supply shock lowers feasible output.
\item Tends to raise prices and lower output.
\end{itemize}

\item \textbf{CAS shifts: labor supply falls}
\begin{itemize}
\item Less labor available $\Rightarrow$ lower full-employment output.
\item Again: $Y\downarrow$, price pressure $\uparrow$.
\end{itemize}

\item \textbf{Demand vs supply shocks: comovement test}
\begin{itemize}
\item Demand shocks: $Y$ and $P$ move together.
\item Supply shocks: $Y$ and $P$ move opposite directions (stagflation).
\end{itemize}

\item \textbf{AD--AS equilibrium (benchmark)}
\begin{itemize}
\item Intersection determines short-run $(Y,P)$.
\item FE is the medium-run reference point.
\end{itemize}

\item \textbf{Demand shock: fiscal expansion}
\begin{itemize}
\item AD right $\Rightarrow$ short run: $Y\uparrow$, $P\uparrow$.
\item Medium run: price adjustment pushes economy back toward FE.
\end{itemize}

\item \textbf{Demand shock: easier monetary policy}
\begin{itemize}
\item AD right via lower rates $\Rightarrow$ $Y\uparrow$, $P\uparrow$ short run.
\item Same comovement as fiscal demand shock.
\end{itemize}

\item \textbf{Supply shock: labor disruption}
\begin{itemize}
\item AS/CAS left $\Rightarrow$ $Y\downarrow$, $P\uparrow$ (stagflation).
\end{itemize}

\item \textbf{Supply shock: productivity declines}
\begin{itemize}
\item AS/CAS left $\Rightarrow$ $Y\downarrow$, $P\uparrow$.
\item Connect to bottlenecks: inflation without a boom.
\end{itemize}

\item \textbf{Summary table: shock $\Rightarrow$ shifts $\Rightarrow$ outcomes}
\begin{itemize}
\item Memorize: which curve shifts and the signs of $(\Delta Y,\Delta P)$.
\end{itemize}

\item \textbf{What this benchmark explains --- and what it cannot}
\begin{itemize}
\item Explains demand vs supply and basic tradeoffs.
\item Cannot handle expectations/persistence well; motivates Phillips/NK.
\end{itemize}

\item \textbf{Preview: sticky prices / NK}
\begin{itemize}
\item Replace ad hoc AS with Phillips curve: inflation depends on expected inflation + slack.
\item Credibility and expectations become central.
\end{itemize}

\item \textbf{Next time}
\begin{itemize}
\item Phillips curve and expectations; then New Keynesian core.
\end{itemize}
\end{enumerate}

\end{document}
